\section{Methodology}
\label{sec:method}

\subsection{Model and Setup}

We use \pythia~\citep{biderman2023pythia}, a 32-layer autoregressive
transformer with $\dmodel = 2560$, accessed via
\transformerlens~\citep{nanda2022transformerlens}. All experiments run on a
single NVIDIA RTX~3090 (24\,GB) with PyTorch~2.4.1. We fix the random seed to
42 throughout. The full experimental pipeline completes in approximately
5~minutes.

\subsection{Concept Direction Extraction}

For each binary concept, we create 20 contrastive prompt pairs---sentences that
differ primarily in the target concept (\eg ``The dog is a \emph{mammal}'' vs.\
``The robin is a \emph{bird}''). We extract residual stream activations at the
last token position across all layers. The concept direction is the normalized
mean difference of positive and negative activations:
\begin{equation}
    \vd_c = \frac{\bar{\vx}^{+}_c - \bar{\vx}^{-}_c}
                 {\|\bar{\vx}^{+}_c - \bar{\vx}^{-}_c\|}
    \label{eq:direction}
\end{equation}
where $\bar{\vx}^{+}_c$ and $\bar{\vx}^{-}_c$ are the mean last-token
activations for the positive and negative prompt sets of concept $c$,
respectively.

\subsection{Concept Taxonomy}

We select 10 binary concepts organized into 5 categories spanning a predicted
composability spectrum:

\begin{enumerate}[leftmargin=*,itemsep=2pt,topsep=2pt]
    \item \textbf{Hierarchical} (expected: compose well): mammal/bird,
    fruit/vegetable---subcategories within a shared ontological hierarchy
    \citep{park2024geometry}.
    \item \textbf{Independent attributes} (expected: compose well):
    positive/negative sentiment, formal/casual register---orthogonal
    properties of text.
    \item \textbf{Semantic relations} (expected: partial): big/small,
    hot/cold---antonym pairs along different perceptual dimensions.
    \item \textbf{Syntactic features} (expected: mixed): past/present tense,
    singular/plural---grammatical properties that may share syntactic
    encoding.
    \item \textbf{Entangled} (expected: compose poorly): happy/sad,
    true/false---concepts whose directions may overlap due to shared
    affective or epistemic dimensions.
\end{enumerate}

From these 10 concepts, we select 10 concept pairs for composition testing,
distributed across categories and spanning the expected composability range.

\subsection{Composition Metrics}

We evaluate compositionality using three complementary metrics, each capturing
a different aspect of what it means for directions to ``compose.''

\para{Metric 1: Orthogonality.}
We compute the cosine similarity between each direction pair:
\begin{equation}
    \cossim(\dA, \dB) = \dA \cdot \dB
    \label{eq:cosine}
\end{equation}
since both directions are unit-normalized. We report both signed cosine
similarity and absolute cosine similarity $\abscossim$, as anti-parallel
directions ($\cossim \approx -1$) interfere as much as parallel directions
($\cossim \approx 1$).

\para{Metric 2: Probing preservation.}
We test whether both concepts remain linearly decodable from the composed
direction $\dAB = \dA + \dB$. For each concept, we train a logistic regression
probe on projections of the original contrastive activations onto $\dAB$ and
measure classification accuracy. We report the \emph{preservation ratio}:
\begin{equation}
    \preservation_c = \frac{\text{Acc}(\text{probe on } \dAB)}
                           {\text{Acc}(\text{probe on } \vd_c)}
    \label{eq:preservation}
\end{equation}
A ratio of 1.0 means no information loss from composition.

\para{Metric 3: Steering composition.}
We add $\alpha \cdot \dAB$ to the residual stream of neutral prompts and
measure logit shifts on concept-relevant tokens. For each component direction,
we compute the \emph{component preservation}: the fraction of the individual
steering effect retained in the composed vector's logit shift. We also measure
\emph{shift alignment}: the cosine similarity between the composed logit shift
and the sum of individual logit shifts.

\subsection{Analysis Protocol}

We analyze all metrics at 6 layers (5, 10, 15, 20, 25, 31), with layer~20 as
the primary analysis layer. Statistical comparisons use one-way ANOVA across
expected composability categories, Spearman rank correlation between
orthogonality and composition metrics, and bootstrap 95\% confidence intervals
(1000 resamples).
